\documentclass[12pt,a4paper,oneside]{article}

\usepackage[utf8]{inputenc}
\usepackage[T1]{fontenc}
\usepackage[brazil]{babel}
\usepackage{geometry}
\geometry{a4paper, margin=2.5cm}
\usepackage{indentfirst}
\usepackage{times}
\usepackage{graphicx}
\usepackage{hyperref}
\usepackage{abntex2cite}

\setlength{\parindent}{1.25cm}
\setlength{\parskip}{0.3cm}

\begin{document}

\begin{center}
  \textbf{\large ESCOLA E FACULDADE DE TECNOLOGIA SENAI GASPAR RICARDO JÚNIOR}\\
  \textbf{Desenvolvimento de Sistemas}\\[1cm]
  \textbf{\Large Lucas Miguel Leite}\\[0.5cm]
  \textbf{Professor: Deivison Takatu}\\[2cm]
  \textbf{\LARGE\textbf{Introdução ao Planejamento Estratégico de TI: Contextualização e Estudo de Caso Itaú Unibanco}}\\[1cm]
  Sorocaba -- 04/02/2026
\end{center}

\vspace{1cm}

\begin{center}
  \textbf{RESUMO}
\end{center}

Este artigo aborda a introdução ao Planejamento Estratégico de TI (PETI), contextualizando a mudança de paradigma da TI de área de suporte operacional para elemento central e estratégico das organizações. A disciplina é apresentada como eixo integrador que conecta fundamentos de TI a Governança, Segurança e Gestão de Serviços. O perfil do profissional de TI é discutido, evidenciando a transição de técnico executor para consultor estratégico que compreende o negócio. Como complemento prático, apresenta-se um estudo de caso do Itaú Unibanco, analisando a transformação digital da instituição, a estratégia de modernização de legado e a parceria com a Amazon Web Services (AWS). Os resultados demonstram que o alinhamento entre TI e estratégias de negócio permitiu ao banco aumentar agilidade, otimizar custos e melhorar a experiência do cliente.

\vspace{0.5cm}
\textbf{Palavras-chave:} Planejamento Estratégico de TI. Transformação Digital. Computação em Nuvem. Itaú Unibanco. AWS.

\newpage

\section{Introdução}

A disciplina de Planejamento Estratégico de TI (PETI) foca na mudança de paradigma da Tecnologia da Informação: de uma área puramente técnica e de suporte para um elemento central e estratégico das organizações. Em um cenário de transformação digital, o PETI é o que viabiliza os objetivos organizacionais, funcionando como eixo integrador que conecta conhecimentos passados (fundamentos de TI) a futuros (Governança, Segurança e Gestão de Serviços) \cite{tarapanoff2015}.

A tríade de desenvolvimento da disciplina estruturase em Conhecimento (visão estratégica aplicada à TI), Habilidade (análise de cenários e avaliação de impactos) e Atitude (pensamento crítico, ética profissional e visão sistêmica). O mercado atual exige a evolução do profissional de TI de técnico executor para consultor estratégico que compreende o negócio, utilizando a TI para gerar inovação, eficiência e vantagem competitiva \cite{rocha2018}.

\section{Organização do Semestre}

O conteúdo é apresentado de forma progressiva em três fases:

\begin{enumerate}
  \item \textbf{Base:} TI e Estratégia Organizacional; Integração TI-Negócio.
  \item \textbf{Aprofundamento:} Tendências tecnológicas, Metodologias de PETI e Impacto nos negócios.
  \item \textbf{Consolidação:} Governança de TI, Segurança da Informação (SGSI), Gestão de Riscos e Serviços.
\end{enumerate}

\section{Estudo de Caso: Planejamento Estratégico de TI do Itaú Unibanco}

A Tecnologia da Informação deixou de ser uma área de suporte operacional para se tornar o núcleo estratégico das instituições financeiras. No contexto atual, a capacidade de processar dados, garantir segurança e oferecer serviços digitais estáveis é o principal diferencial competitivo dos bancos \cite{porter1989}.

O Itaú Unibanco, maior instituição financeira privada da América Latina, enfrentou o desafio de competir com fintechs que possuem estruturas de custo menores e tecnologias mais ágeis. Para manter sua liderança, o banco reformulou seu PETI, focando na descentralização de sistemas legados e na adoção massiva de computação em nuvem.

\subsection{Contexto e Desafios}

Até meados de 2015, a arquitetura tecnológica dos grandes bancos brasileiros era predominantemente baseada em mainframes monolíticos. Embora robustos e seguros, esses sistemas apresentavam dificuldades para integração com novas aplicações móveis e limitavam a velocidade de inovação. O Itaú identificou que a manutenção dessa estrutura gerava alto custo de data centers físicos, lentidão no desenvolvimento de novos produtos (time-to-market elevado) e dificuldade em escalar serviços em momentos de pico \cite{itauchold2025}.

\subsection{Estratégia de Modernização e Parceria com AWS}

Em 2020, o banco firmou parceria de longo prazo com a Amazon Web Services para migrar a maior parte de sua infraestrutura de TI. A estratégia baseou-se em três pilares \cite{itauchold2020}:

\begin{enumerate}
  \item \textbf{Modernização de Aplicações:} Reescrita de sistemas legados utilizando arquitetura de microsserviços, permitindo que falhas em um serviço não derrubassem todo o sistema bancário.
  \item \textbf{Cultura de Dados e IA:} Uso intensivo de Analytics e Inteligência Artificial para personalizar a oferta de crédito e investimentos em tempo real.
  \item \textbf{Mudança Cultural (Agile):} Adoção de metodologias ágeis e organização em squads, unindo profissionais de TI e Negócios no mesmo time.
\end{enumerate}

\subsection{Resultados e Impactos no Negócio}

A modernização permitiu ao Itaú suportar o crescimento exponencial das transações digitais, impulsionado pela criação do Pix. Os resultados incluem redução do Custo Total de Propriedade (TCO) com desativação progressiva de mainframes, aumento de eficiência (processos que levavam horas passaram a minutos na nuvem) e estabilidade com escalabilidade automática durante picos de acesso \cite{itauchold2023}.

\section{Conclusão}

A análise do caso Itaú Unibanco demonstra a aplicação prática dos conceitos de PETI. A tecnologia não foi tratada como fim em si mesma, mas como meio necessário para garantir a sobrevivência e o crescimento em um mercado disruptivo. O sucesso deve-se não apenas à contratação de fornecedores como a AWS, mas à reestruturação da cultura organizacional e ao alinhamento profundo entre os objetivos de negócio e as diretrizes de TI. Este caso reforça que o PETI deve ser dinâmico, integrado à alta gestão e focado em gerar valor perceptível \cite{kluyver2007}.

\bibliographystyle{plain}
\bibliography{referencias}

\end{document}
